\section{建康、南渡与一座受限的朝廷}
\label{sec:eastern-jin-establishment-southern-society}

永嘉元年，司马睿出镇建业。对于正在北方内战中失去支点的西晋宗室而言，这不是把一套完备朝廷搬到长江以南，而是先在一处既有的区域中心寻找可以接续诏令、官员与军队的关系网。王导等人参与联络江东士族、整顿东南军政，遂使司马睿的府署逐渐成为南方政治中心的雏形。后来以建康为名的这座城市，因而并非只是一座避难的城池：它同时是晋室、先到的官僚与当地精英必须共同使用、又各自试图影响的地点。

《晋书·元帝纪》与〈王导传〉保存了司马睿出镇和王导参与的基本线索，《建康实录》则提供都城视角下的地名与人物线索。三者都不足以把最初的结盟写成一次众望所归的迎接。谁在何时表示支持、府署究竟能够调动多少州郡，以及江东旧族如何看待北来的宗室，传世文字并没有同样清楚的答案。所谓“侨姓”与“江东旧族”也不是两个内部毫无差别的整体；把后来形成的门第格局倒投到三〇七年，只会掩盖建业政权最初的试探与协商。

\subsection{城破、渡江与可用资源的重新组合}

北方的灾难使这种试探变得急迫。三一一年洛阳陷落，怀帝被俘；三一三年怀帝被杀，长安另立的愍帝朝廷也在三一六年随着刘曜攻陷长安而终结。几次都城的失守并不自动规定每一个人的去向，却使原有的官僚任命、军队供给和家族居处失去共同的中心。南下者带来的不只是人口，也有官职资格、旧有从属关系、土地诉求和对北方局势各不相同的消息。建业的府署能够借此扩大，却也必须面对安置、征发和地方秩序的压力。

因此，“南渡”不应被想象成整齐的一次迁徙，更不能由正史中带有修辞色彩的死散、流亡数字推算出精确的人口总量。有人抵达江东，有人留在淮北、河南或关中，也有人在行进、停留和再迁徙之间改换依附。侨置州郡和侨居名目的出现，说明朝廷试图在籍贯记忆、任官资格与实际居住之间建立行政秩序；它们却不能单独证明某一地已经聚居了多少人，或当地社会已经接受了何种安排。南方朝廷的社会基础，正是在这种不均匀的流动、交涉与征发中形成，而不是由一个抽象的“北来士族”一次建立。

建兴元年祖逖率部北渡，在豫州、河南一带屯田并与地方力量结盟，显示建业一方并未放弃江北联系。可这类经营依赖将领本人的人际网络、能够获得的粮食与当地关系，并不等于一条由建业稳定供给、统一指挥的北伐战线。《晋书·祖逖传》与《资治通鉴》足以确认北渡、屯田和经营的轮廓；它们所载收复范围、兵粮数额和归附情形却不宜照录为精确地图。《世说新语》注所引《晋阳秋》可补充后世流传的线索，仍不能把人物声名替代为军政控制的证据。

\subsection{从晋王到皇帝}

长安陷落之后，三一七年司马睿在建康即晋王位，南方朝廷由区域府署转而公开承担晋室继承的名义。次年称帝、改元，东晋才以建康为都正式确立。帝号解决的是“谁可以代表晋”的迫切问题，并没有同时解决“谁能征发兵粮、约束军镇、处理流民”的问题。北方多政权并立，长江、淮河与江汉之间的道路又分别联系着不同的军镇和地方社会；建康所能直接触及的资源，远小于西晋全盛时洛阳所宣称的天下。

《晋书·元帝纪》、〈王导传〉与《资治通鉴》对即王、称帝的前后次序可以互相校核，《建康实录》也有助于辨认都城叙事的线索。不过，受禅礼仪究竟由哪些北方士人代表，南方地方社会以何种方式接受新皇统，均不能从这些以朝廷人物为中心的材料中直接得出。把三一八年写成旧秩序在江南无损地续存，或写成一个完全脱离北方的新国家，同样会失真。东晋首先是一个以晋朝名义重新组织的南方朝廷；它的财政、兵源和北方目标，始终受区域分裂所限制。

\subsection{平乱的军镇与王敦的压力}

这种限制很快显露在长江中游。大兴三年，荆湘地区的杜弢叛乱被王敦等人平定，朝廷重新取得该地区一部分军政控制。战争的结果不能简单化为“秩序恢复”：永嘉以后的人口流动、州郡征发与地方武装的竞争，都可能使不同群体卷入冲突。《晋书·杜弢传》与〈王敦传〉保留的是以朝廷和将领为中心的叙述，难以据此逐一判定流民、地方豪强及史书所称其他群体的参与程度。较为清楚的是，建康要维系长江中游，必须依靠能够自行调兵、筹给的军镇将领；平乱也因此增加了王敦的军事声望。

大兴四年，王敦以“清君侧”为名自武昌东下，迫使元帝调整中枢人事，并承认其军事实力。这不是一个只由个人好恶引起的宫廷插曲。东晋初建时，王氏的政治网络和荆州方面的兵力原是朝廷得以运转的条件；当元帝试图扶植近臣、重建较可控制的中枢时，外镇便把对任用和诏令的争议转化为武力压力。《晋书·元帝纪》、〈王敦传〉和〈王导传〉对各方的言辞、调停与责任分配并不完全一致，《资治通鉴》的连贯叙述也不能消除这些差异。可以把握的是，建康不得不让步，而皇帝对江上军镇的约束由此显著削弱。

永昌元年，王敦再度起兵，攻入建康，杀逐元帝近臣并控制朝政。城内战斗的经过、元帝每一步的应对和死者名单，在《晋书》诸传之间仍有需要核对之处；不必用戏剧化的宫门细节填补材料的空白。事件本身已经说明，建康虽是皇统所在地，却不能仅凭都城、礼仪和官号支配沿江的军事资源。王敦没有立即废帝，反而更能显示当时的结构：晋朝名义仍有价值，能够决定日常朝局的力量却可以来自都城之外。

\subsection{史料所照见的朝廷与未被照见的社会}

这一段年序以《晋书》的帝纪、列传为主干。它成书于唐初，汇集较早史料，保存了司马睿、王导、祖逖与王敦行动的传世叙事；其人物褒贬、事后归责和宫廷秘闻却不能直接当作三、四世纪的现场记录。《资治通鉴》将这些事迹编入连续年代，特别适合核看都城陷落、即位与起兵的前后关联，但它是北宋的重新编排，不是独立的建康档案。《建康实录》保留都城和地方线索，也因成书较晚、传本与引据复杂，不能替代前两者。

城址、墓葬、墓志以及侨州郡的地方性记述，能够从另一种尺度提示城市空间、仕宦网络、籍贯自述和迁徙留下的痕迹；它们同样有发掘范围、墓主自我呈现、追赠官衔和传本层次的限制。一处墓葬不能代表整座建康，一条侨置记载也不能变成人口普查。本节因而只把它们作为校正正史视野的必要材料类型，而没有以考古或地方记录虚构三一七年拥立的现场、南渡人口的总数，或江东社会一致的态度。可确证的，是朝廷在建康重新竖起了晋室名号，并在迁徙、地方军镇和北方分裂的牵制下求取存续；其余那些无法由现存材料逐项回答的社会经验，应当保留为问题，而非用后来的整齐叙事加以填平。
